\chapter{MicroPython IDE} \label{ide}

\section{Connecting Your Microcontroller}
An \textbf{IDE} (Integrated Development Environment) is a program that puts everything you need
for coding in one place: a file browser, an editor, and buttons to run and stop your code.
Throughout this guide, you will use a free online IDE to connect to the microcontroller and to
open, edit, and run the code stored on it. You can access the IDE from your browser by going to
\url{https://viper-ide.org/}.

\begin{tcolorbox}[colback=yellow!10!white,colframe=yellow!50!black]
    \textbf{Use Chrome or Edge.} The IDE talks to the microcontroller over USB using a browser
    feature called Web Serial, which only Chrome, Edge, and other Chromium-based browsers support.
    It will not connect from Safari, Firefox, or an iPad.
\end{tcolorbox}

Click on the USB icon in the top right and choose your microcontroller from the list:

\begin{figure}[H]
    \centering
    \maybeincludegraphics[width=.6\linewidth]{common/viper_ide_usb_connect.png}
    \caption{Click the button highlighted in red.}
\end{figure}

If you see multiple items in the dialog that pops up, choose the one that starts with "USB JTAG". See below for an example:
\begin{figure}[H]
    \centering
    \maybeincludegraphics[width=.45\linewidth]{common/viper_ide_usb_choice_connect.png}
    \caption{Click the button highlighted in red.}
\end{figure}

Once you have connected, you will see a green dialog labeled "Device connected", and the file
manager on the left will fill in with the files stored on the microcontroller. All of the project
files are already there.
\begin{figure}[H]
    \centering
    \maybeincludegraphics[width=\linewidth]{common/viper_connected.png}
    \caption{Make sure to open the right file for this project}
\end{figure}

\section{The Serial Console}
Below the editor is a panel labeled \textbf{Terminal}. This guide calls it the
\textbf{serial console}. Whenever a program calls \texttt{print()}, the
text travels over the USB cable from the microcontroller and appears here. Every project uses
the serial console to report what it is doing, so it is the first place to look when something
is not working.

After your device is connected, refer back to the chapter you are working on for the next steps.
